\subsection{Expert-selection comparison with SteerMoE}
\label{sec:expert-selection-audit}

The comparison separates expert detection from the intervention applied to
those experts. SteerMoE uses paired contrastive inputs and counts top-$k$
expert inclusions separately in the two populations. Its score is the
token-pooled activation-rate difference
\begin{equation}
\Delta_{\ell e}^{\mathrm{paper}} =
\frac{A_{\ell e}^{+}}{N^{+}}-\frac{A_{\ell e}^{-}}{N^{-}},
\end{equation}
where $A$ counts tokens selecting the layer--expert identity and $N$ counts
all tokens in the corresponding population. Detection therefore uses routing
frequencies, not hidden-state similarity. Positive and negative differences
identify opposite behavioral associations; steering changes the selected
experts' log-softmax scores to the original maximum plus $\epsilon$ or minimum
minus $\epsilon$~\cite{fayyaz2025steermoe}, Sections 3.1--3.2.
The authors' released custom-steering notebook ranks layer--expert identities
together by absolute risk difference\footnote{\url{https://github.com/adobe-research/SteerMoE/blob/main/custom_steering.ipynb}}.
Our policy instead assigns a separate budget to each of a preselected set
of layers.

\paragraph{Implemented estimator.}
For calibration attempt $r$, let $y_r\in\{0,1\}$ be correctness,
$f_{r\ell e}$ the fraction of saved routing tokens selecting expert $e$ at
layer $\ell$, and $w_r=1/m_{p(r)}$, where $m_{p(r)}$ is the number of usable
attempts for its source problem. With $W=\sum_r w_r$, the implemented score is
\begin{equation}
\pi=\frac{\sum_r w_r y_r}{W},\qquad
L_{\ell e}=\frac{\sum_r w_r f_{r\ell e}y_r}
                       {\sum_r w_r f_{r\ell e}}-\pi.
\end{equation}
Thus calibration gives every problem total weight one and first normalizes
routing counts within each trace. It contrasts observed correct and incorrect
completions; it does not construct paired counterfactual behavioral inputs.
Experts need at least four supporting problems in the audited policies.
Within each selected layer, eligible experts are ordered by signed lift,
then support, then expert ID, and capped at native top-$k$
(eight for Qwen and four for GPT-OSS). Fixed bias scales each
retained lift by the largest retained magnitude; paper-style steering uses
only membership in the selected set, not those relative magnitudes.

\paragraph{Why activation-based scores need not select the same experts.}
Define $\mu_1=E_w[f_{r\ell e}\mid y_r=1]$,
$\mu_0=E_w[f_{r\ell e}\mid y_r=0]$, and
$\bar f=E_w[f_{r\ell e}]$. Algebra applied to our estimator gives
\begin{equation}
L_{\ell e}=\frac{\pi(1-\pi)}{\bar f}(\mu_1-\mu_0),\qquad \bar f>0.
\end{equation}
Under these same weights, lift and activation-rate difference have the same
sign but can have different rankings: lift additionally divides by the
expert's mean activation frequency. This identity does not equate either
score to the paper's token-pooled estimate on different, paired inputs.

The saved frequency masses, activation-weighted accuracies, and calibration
baseline suffice to reconstruct $\mu_1-\mu_0$ without replaying a model.
Table~\ref{tab:selection-audit} changes only the ranking score and retains
our weights, sign filter, support threshold, layers, and target budget.
The audit also reconstructs every saved selected set and normalized score.
Only 14 of 48 Qwen positive layer--expert identities remain selected;
the four-target GPT-OSS policies retain 6 of 16 positive identities and
7 of 16 negative identities.
This demonstrates selection sensitivity even before changing the data
construction or token weighting. It does not measure intervention performance
for these alternative sets.

\begin{table}[ht]\centering\small
\caption{Selection sensitivity using saved calibration sufficient statistics. RD here uses the same problem-balanced trace weights, support filter, polarity, layer set, and per-layer budget as lift; it is not a reproduction of the paper\textquotesingle s token-pooled contrastive estimator. Overlap counts layer--expert identities.}
\label{tab:selection-audit}
\begin{tabular}{llrrl}\toprule
Model & Polarity & Layers & Overlap & Overlap per layer \\\midrule
GPT-OSS & negative & 4 & 7/16 & 15: 1/4, 19: 2/4, 21: 1/4, 23: 3/4 \\
GPT-OSS & positive & 4 & 6/16 & 15: 2/4, 19: 1/4, 21: 1/4, 23: 2/4 \\
Qwen & positive & 6 & 14/48 & 27: 4/8, 34: 3/8, 35: 2/8, 36: 1/8, 38: 3/8, 39: 1/8 \\
\bottomrule\end{tabular}\end{table}


\paragraph{Interpretation of the intervention comparison.}
The completed Qwen fixed runs share selection statistics and target sets.
The completed GPT-OSS fixed and paper-style conditions use four targets
per layer, with the same selected identities within each polarity.
Their results compare intervention rules using our selector; they do not
evaluate the full SteerMoE detection and steering pipeline.

The completed OSS comparison already holds the selected identities fixed
within each polarity while changing the intervention rule. A selector
comparison should instead fix the split, generation contract, layers, budget,
and intervention, changing only the ranking estimator. The present OSS
budget is four targets per layer, matching native dispatch; smaller budgets
of one or two targets would test less extensive promotion. Matched random
sets would help separate selector information from generic perturbation. A faithful detection replication additionally requires
paired contrastive inputs and token-pooled activation counts. The present
audit changes no policies and launches no generation runs.
